\subsection{序群的初等运算}

设H是序群G的一个子群；显然，G的序在H上的限制与H的群结构相容；除非另有说明，我们总按这种方式给H赋序。若P是G的正元素集合，则H的正元素集合是 \(H \cap P\) 。

设 \((\mathbf{G}_{\alpha})\) 是一族序群；根据有序集乘积的定义（集合论，III，第137页），乘积群 \(\mathbf{G} = \prod_{\alpha}\mathbf{G}_{\alpha}\) 装备着

一个序关系，G 的两个元素之间的关系 \(\ll (x_{\alpha}) \leqslant (y_{\alpha})\) 按定义等同于“ \(\ll x_{\alpha} \leqslant y_{\alpha}\) 对所有 \(\alpha\) ”。显然，这个序关系与G的群结构相容；它使G成为一个序群，我们称之为序群 \(G_{\alpha}\) 的乘积。G的正元素是那些所有分量均为正的元素。当所有因子 \(G_{\alpha}\) 都等于同一个序群H时，G就是由指标集I到H的映射组成的群 \(H^{I}\) ，从I到H的两个映射之间的关系 \(\ll f \leqslant g\) 就是“ \(\ll f(\alpha) \leqslant g(\alpha)\) 对所有 \(\alpha \in I\) ”；正映射是指只取正值的映射。序群族 \((G_{\alpha})\) 的直和定义为其乘积的一个序子群（II，第202页）。

设 \((\mathbf{G}_{\iota})_{\iota \in I}\) 是一族序群，其指标集 I 是良序的；回忆（集合论，III，第157页），在乘积集 \(\mathbf{G} = \prod_{\iota} \mathbf{G}_{\iota}\) 上定义了一个称为字典序的序关系，关系 \(\ll (x_{\iota}) < (y_{\iota})\)

（G 的两个元素之间）按定义等同于“若 β 是使得 \(x_{i} \neq y_{i}\) 的指标 i 中最小者，则 \(x_{\beta} < y_{\beta}\) ”。回忆，全序集的良序族的乘积在字典序下仍是全序的。在一般情形下，直接验证可知G上的字典序与其群结构相容；装备此序后，群G因此是一个序群，称为良序序群族的字典序积（ \(G_{i}\) ）。

注记. --- 1) 在最常见的情形中，良序指标集 I 将是一个有限区间 \([1, n]\) （包含在 \(\mathbf{N}\) 中）。

\begin{enumerate}
\def\labelenumi{\arabic{enumi})}
\setcounter{enumi}{1}
\tightlist
\item
  字典序乘积G的正元素集合由0以及那些非零元素组成，这些非零元素的最小索引非零分量为正。
\end{enumerate}

